\section{OBJECTIVE AND SOURCE REVIEW}
\textbf{Erd\H{o}s \#62; independent attempt, 2026-09-23.}
Must two graphs of chromatic number $\aleph_1$ share a four-chromatic subgraph, or even one of chromatic number $\aleph_0$? I read the entire \texttt{62.tex}. Its necessary odd-girth and size conditions concern finite critical-subgraph spectra. Here I check what information such spectra can retain about uncountable chromatic number.
\section{WORK: A COUNTABLE SPECTRUM MODEL}
For any graph $G$ of infinite chromatic number, choose one finite graph $F_i$ from each isomorphism type of finite subgraph of $G$, and form their disjoint union $H=\bigsqcup_i F_i$. There are countably many finite graphs, so $H$ is countable. It has chromatic number exactly $\aleph_0$: if all $F_i$ were $r$-colourable for some fixed finite $r$, compactness would give an $r$-colouring of $G$, contrary to hypothesis. Compactness here is the finite intersection property in the product $\{1,\ldots,r\}^{V(G)}$, with the closed edge constraints.
Every finite connected graph $F$ embeds in $H$ if and only if it embeds in $G$. The forward direction holds because the image of a connected graph lies in one component $F_i$; the reverse direction holds because a representative of $F$ is one of the components. In particular every finite vertex-critical graph is connected: otherwise a component of maximum chromatic number survives deletion of a vertex from another component. Thus
\[
\mathcal C_4(H)=\mathcal C_4(G).
\]
The same construction works for every finite critical chromatic level.
\section{ADVERSARIAL VERIFICATION AND REMAINING GAP}
Disconnected finite spectra need not be identical, since the union may create new disjoint configurations; connectedness is why the critical-spectrum assertion is correct. A countable graph can never have chromatic number $\aleph_1$. Therefore realizing a desired finite spectrum by a disjoint union, even one with unbounded finite chromatic number, cannot supply the uncountably chromatic counterexample required by the problem. Conversely this observation does not prove that any prescribed countable spectrum has an uncountably chromatic realization. That realization obstacle remains untouched. No general common four-chromatic graph is found, and the stronger countably chromatic variant also remains unresolved.
\section{SOURCES CHECKED}
Complete local source: \texttt{gpt\_pro\_5.4/62.tex}. Official indexed formulation checked at \url{https://www.erdosproblems.com/62}, 2026-09-23. The countable-model observation is elementary and is not claimed as a novel graph theorem.
\section{FINAL}
\textbf{UNRESOLVED.} Finite critical spectra admit countable models, exposing a specific limit of finite-spectrum constructions.
\section{COMPLETION ESTIMATE}
Conservative subjective progress toward both common-subgraph questions.
\textbf{COMPLETION: 5\%}
